% !PRINTED 353
\weberpage{353}

The last equation yields the distinction between the two cases in which the equation c) has two positive roots or none.

For by the rules of the differential calculus one obtains for the maximum of $\sin\Theta^{2m}\cos\Theta^{2n}$
\[
\frac{m^m n^n}{(m+n)^{m+n}},
\]
which is attained for $\mathrm{tg}\,\Theta^2 = m:n$. If therefore $\lambda$ lies below this bound, we have two real roots in c), otherwise none. In the formulae (2), $\lambda$ is a given quantity, and one now has to seek from the tables the angle $\Theta$ that satisfies these equations. If one as yet has no knowledge at all of the approximate position of this angle, it is expedient, for the first approximation, to use a table proceeding say from degree to degree and abridged to two decimal places, in order then to improve the value so found with the help of more exact tables.

\gesperrt{Gauss} uses not the trigonometric tables, but the addition and subtraction logarithms first introduced by him. We shall show this briefly in one of the cases. The details for the practical application of the method are to be sought in \gesperrt{Gauss}'s memoir $^1$).

The tables of addition and subtraction logarithms, as they were first introduced and computed by \gesperrt{Gauss}, and as they are now also to be found in the tables in common use, give, for three numbers which are greater than 1:
\[
a, \quad b = 1+\frac{1}{a}, \quad c = 1+a,
\]
the \gesperrt{Briggian} logarithms:
\[
A, \quad B, \quad C.
\]
If $\Theta$ is an angle in the \gesperrt{first} octant, one may set:
\[
(3)\quad a = \mathrm{cotg}\,\Theta^2, \quad b = \frac{1}{\cos\Theta^2}, \quad c = \frac{1}{\sin\Theta^2}, \quad 0 < \Theta < 45^0,
\]
and if $\Theta$ lies in the \gesperrt{second} octant:
\[
(4)\quad a = \mathrm{tg}\,\Theta^2, \quad b = \frac{1}{\sin\Theta^2}, \quad c = \frac{1}{\cos\Theta^2}, \quad 45^0 < \Theta < 90^0.
\]

\noindent\rule{4cm}{0.4pt}

$^1$) Beiträge zur Theorie der algebraischen Gleichungen, second part (1849). \gesperrt{Gauss} Werke, Vol. III, p. 85.
